\justifying
\NSFCBodyText

% \subsubsection{已具备的实验条件}

% \subsubsubsection{硬件设施}
% 依托单位文化产业研究中心配备完善科研设施：（1）高性能计算集群：NVIDIA A100 GPU 4张（每张40GB显存），总内存512GB；（2）大容量数据存储系统：总容量500TB，支持多用户并发访问和自动备份；（3）专业级影像处理工作站：高色域显示器（99\% Adobe RGB）和专业级图形卡；（4）专业录音与视频编辑设备。表\ref{tab:hardware_config}详细列出了主要硬件设施配置。

% \begin{table}[htbp]
%   \centering
%   \fontsize{11}{11}\selectfont
%   \caption{\textbf{主要硬件设施配置}}
%   \begin{tabular}{lccc}
%   \toprule
%   \textbf{设备类型} & \textbf{配置规格} & \textbf{数量} & \textbf{主要用途} \\
%   \midrule
%   GPU计算卡 & NVIDIA A100 40GB & 4张 & 深度学习训练 \\
%   服务器内存 & DDR4 ECC & 512GB & 大规模数据处理 \\
%   存储系统 & 企业级NAS & 500TB & 多模态数据存储 \\
%   工作站 & 高色域显示器 & 5台 & 影像处理与分析 \\
%   录音设备 & 专业级麦克风 & 2套 & 访谈录制 \\
%   \bottomrule
%   \end{tabular}
%   \captionsetup{font={footnotesize,stretch=1.25},justification=raggedright}
%   \label{tab:hardware_config}
% \end{table}

% \subsubsubsection{软件资源}
% 实验室部署完整数据分析工具链：（1）自然语言处理：NLTK、spaCy、jieba、HanLP等；（2）深度学习框架：TensorFlow、PyTorch、Keras，支持BERT、GPT等预训练模型；（3）计算机视觉：OpenCV、PIL、scikit-image、detectron2；（4）数据可视化：Tableau、D3.js、Matplotlib、Seaborn；（5）社会科学统计：SPSS、Stata、R；（6）项目管理：Git、GitHub、Slack、Trello。

% \subsubsubsection{数据库资源}
% 依托单位图书馆购买丰富学术数据库：（1）外文数据库：Web of Science、Scopus、JSTOR、SpringerLink、ScienceDirect、EBSCO、SAGE、Wiley等；（2）中文数据库：中国知网（CNKI）、万方数据、维普资讯等；（3）专业数据库：MyDramalist、AsianWiki、Oricon等；（4）统计年鉴：中国统计年鉴、日本统计年鉴等。

% \subsubsection{尚缺少的实验条件及拟解决途径}

% \subsubsubsection{缺少的实验条件}
% （1）日本本土社交媒体数据因地域限制和平台政策无法完整获取；（2）日本娱乐产业一手资料属商业机密；（3）韩国、欧美娱乐产业数据资源分散；（4）大规模视频数据存储与处理需更专业设备；（5）日语情感分析缺乏高质量标注语料。

% \subsubsubsection{拟解决途径}
% （1）国际合作：与东京大学信息学环、早稻田大学文学学术院、首尔国立大学言论情报学研究所、阿姆斯特丹大学传媒系、南加州大学新闻与传播学院等建立合作关系；（2）数据采购：通过 Brandwatch、Talkwalker、Cision 等获取社交媒体历史数据（预算3-5万元/年）；（3）公开数据集：利用 Twitter Academic Research Product Track 等；（4）自主研发：已开发网络爬虫工具和多语言情感分析工具原型；（5）语料库建设：与日本合作机构联合建设日语情感标注语料库（约10万条）；（6）设备升级：申请经费升级视频存储和处理设备（预算5-8万元）。

% \subsubsection{利用国家重点实验室等研究基地的计划}

% \subsubsubsection{合作单位}
% 申请团队计划与国内高水平研究基地建立紧密合作关系：中国人民大学新闻学院社会发展研究基地、清华大学新闻与传播学院新媒体研究中心、中国传媒大学文化产业管理学院、文化和旅游部文化产业研究中心、国家广播电视总局广播电视规划院。

% \subsubsubsection{合作方式}
% 学者互访（每年邀请合作基地知名学者来校讲学1-2次，项目负责人赴合作基地交流访问1-2次）；联合研究（在重大理论问题上开展联合研究，共同发表高水平论文）；资源共享（共享部分研究数据和设备资源）；人才培养（共同指导研究生，开展学生访学交流）。


\subsubsection{已具备的实验条件}

本项目依托申请人所在单位的高水平科研平台，在硬件设备、算力资源、数据积累及团队梯队等方面已具备了坚实的研究基础，能够完全保障项目的顺利实施。

\paragraph{\textbf{1. 硬件平台与机器人实验环境}}
实验室与国内具身智能领军企业（\textbf{星辰智能}）建立了深度科研合作伙伴关系，已在实验室部署了**履带式双臂人形机器人**作为本项目的主力验证平台。该机器人配备高精度灵巧手、多自由度头颈模组及双目深度相机，开放了底层的关节力矩控制接口（SDK）与运动学解算库，完全满足Social-MLA模型从算法验证到真机部署的硬件需求。此外，实验室配备了**OptiTrack光学动作捕捉系统**与**Xsens惯性动捕设备**，能够实时采集高精度的真实人体社交动作数据，用于算法的基准真值（Ground Truth）对比与评估。

\paragraph{\textbf{2. 高性能计算与仿真算力资源}}
申请人所在研究院计算中心已建成大规模高性能计算集群，配备**NVIDIA A100/H100 Tensor Core GPU**节点共计（\textbf{填写数量}）个，存储容量达PB级，并部署了InfiniBand高速互联网络。该算力平台已预装NVIDIA Isaac Lab、Omniverse等物理仿真引擎及PyTorch深度学习框架，能够充分支撑Social-MLA大模型的分布式训练以及百万级并发的Sim-to-Real强化学习仿真任务。

\paragraph{\textbf{3. 前期数据与算法代码库}}
项目组前期已积累了约（\textbf{200小时以上}）的高质量互联网演讲与交互视频数据，并完成了初步清洗与SMPL-X参数化处理，构建了\textbf{Embodied-Social-Pre}预训练数据集。同时，团队自主研发了包含Motion Retargeting、语音驱动表情生成及全身协同生成的全栈代码库，为本项目的快速启动奠定了软件基础。

\subsubsection{利用国家/重点实验室等研究基地的计划与落实情况}

本项目将紧密依托（\textbf{填写依托单位的国家重点实验室或部门重点实验室名称，例如：虚拟现实技术与系统全国重点实验室 / 机器人学国家重点实验室}）开展研究工作，计划落实情况如下：

\paragraph{\textbf{1. 依托重点实验室的理论指导与学术交流}}
项目将纳入重点实验室的年度研究计划体系，接受（\textbf{填写院士/实验室主任姓名}）院士及学术委员会的定期指导。充分利用实验室在``人机共融''与``虚拟现实''领域的顶级学术网络，定期邀请国内外具身智能领域的专家进行研讨，确保项目技术路线始终处于国际前沿。

\paragraph{\textbf{2. 共用重点实验室的特种实验设施}}
计划利用重点实验室的**大型沉浸式交互实验场**（\textbf{或填写具体的实验室设施名称}），开展人机共身交互的主观评测实验。该实验场具备可控的光照、声场环境及全方位生理信号监测设备（如脑电EEG、皮电EDA），为评估老年人在与机器人交互过程中的真实情绪反馈与认知负荷提供了科学的实验环境。

\subsubsection{尚缺少的实验条件和拟解决的途径}

尽管项目组已具备核心研发条件，但在**真实适老化场景的复杂性验证**及**极端工况下的数据完备性**方面仍存在一定缺口，拟通过以下途径解决：

\paragraph{\textbf{1. 尚缺少：真实居家/养老机构的复杂交互场景}}
\textbf{拟解决途径：} 依托合作企业（\textbf{公司名}）的行业渠道资源，计划在项目第三年度，与（\textbf{填写具体的养老院/社区名称，或写“北京市某主要养老示范社区”}）建立**``具身智能适老化应用示范点''**。将实验从实验室的结构化环境延伸至真实的非结构化生活场景，在严格的安全监管下收集真实用户的交互反馈数据，以验证系统的鲁棒性与社会接受度。

\paragraph{\textbf{2. 尚缺少：机器人极端工况下的本体感知数据}}
\textbf{拟解决途径：} 真实机器人难以进行破坏性实验（如剧烈碰撞、跌倒），导致极限边界数据缺失。项目组将充分利用**Sim-to-Real技术**，在Isaac Lab仿真环境中构建高保真的``数字孪生''场景，通过程序化生成（Procedural Generation）技术合成大规模的极端交互数据（如拥挤人群、复杂障碍物），通过**合成数据增强（Synthetic Data Augmentation）**的方式填补真实数据的空白，提升策略网络的泛化边界。